% Part: second-order-logic
% Chapter: sol-and-sets
% Section: comparing-sets

\documentclass[../../../include/open-logic-section]{subfiles}

\begin{document}

\olfileid{sol}{set}{cmp}

\olsection{సమితులను పోల్చడం}

\begin{prop}
\tetoken{సూత్రం}{!!{formula}} $\lforall[x][(X(x) \lif Y(x))]$ ఉపసమితి
సంబంధాన్ని నిర్వచిస్తుంది; అంటే
$\Sat{M}{\lforall[x][(X(x) \lif Y(x))]}[s]$ అయితేనే $s(X)
\subseteq s(Y)$.
\end{prop}

\begin{prop}
\tetoken{సూత్రం}{!!{formula}} $\lforall[x][(X(x) \liff Y(x))]$ సమితులపై
తాదాత్మ్య సంబంధాన్ని నిర్వచిస్తుంది; అంటే
$\Sat{M}{\lforall[x][(X(x) \liff Y(x))]}[s]$ అయితేనే
$s(X) = s(Y)$.
\end{prop}

\begin{prop}
\tetoken{సూత్రం}{!!{formula}} $\lexists[x][X(x)]$ శూన్యేతరం అనే ధర్మాన్ని
నిర్వచిస్తుంది; అంటే $\Sat{M}{\lexists[x][X(x)]}[s]$ అయితేనే
$s(X) \neq \emptyset$.
\end{prop}

$X$ సమితి $Y$ సమితిని పరిమాణంలో మించకపోవడం,
$\cardle{X}{Y}$, అంటే \tetoken{ఒకటి-ఒకటి ప్రమేయం}{!!a{injective}}
$f\colon X \to Y$ ఉండటం. ఒక ప్రమేయం ఒకటి-ఒకటి అని, అలాగే~$X$లోని
ప్రవేశాలకు దాని విలువలు~$Y$లో ఉన్నాయని వ్యక్తం చేయగలం కాబట్టి,
వ్యక్తి క్షేత్రపు ఉపసమితులపై పరిమాణంలో మించని సంబంధాన్ని కూడా
నిర్వచించవచ్చు.

\begin{prop}
ఈ సూత్రం
\[
\lexists[u][(\lforall[x][(X(x) \lif Y(u(x)))] \land \lforall[x][\lforall[y][(\eq[u(x)][u(y)] \lif
      \eq[x][y])]])]
\]
పరిమాణంలో మించని సంబంధాన్ని నిర్వచిస్తుంది.
\end{prop}

రెండు సమితులు ఒకే పరిమాణం గలవి, లేదా “సమసంఖ్యాకాలు,”
$\cardeq{X}{Y}$, అంటే \tetoken{ఒక ద్విగుణ ప్రమేయం}{!!a{bijective}}
$f\colon X \to Y$ ఉండటం.

\begin{prop}
ఈ సూత్రం
\begin{multline*}
  \lexists[u][(\lforall[x][(X(x) \lif Y(u(x)))] \land {}\\
    \lforall[x][\lforall[y][((X(x) \land X(y) \land
      \eq[u(x)][u(y)]) \lif \eq[x][y])]]
  \land {} \\\lforall[y][(Y(y) \lif \lexists[x][(X(x)
      \land \eq[y][u(x)])])])]
\end{multline*}
సమసంఖ్యాక సంబంధాన్ని నిర్వచిస్తుంది.
\sourcecorrection{OLTESOLSET-006}{ఆంగ్ల మూలంలోని సమసంఖ్యాకత సూత్రం $u$ను మొత్తం వ్యక్తి క్షేత్రంపై ఒకటి-ఒకటిగా కోరింది. $X$, $Y$లకు ఒకే పరిమాణం ఉన్నా వాటి పూరకాలకు వేర్వేరు పరిమాణాలు ఉండవచ్చు కాబట్టి, $X$ నుంచి $Y$కు ఉన్న ద్విగుణ ప్రమేయం మొత్తం వ్యక్తి క్షేత్రంపై ఒకటి-ఒకటి ప్రమేయంగా విస్తరించకపోవచ్చు. ఒకటి-ఒకటితన నిబంధనను $X$ ప్రవేశాలకు పరిమితం చేశాం; మిగతా రెండు నిబంధనలు $u$ను $X$ నుంచి $Y$పైకి ద్విగుణ ప్రమేయంగా చేస్తాయి.}
\end{prop}

ఈ \tetoken{సూత్రాలను}{!!{formula}s} వరుసగా $X \subseteq Y$,
$X = Y$, $X \neq \emptyset$, $\cardle{X}{Y}$, ఇంకా
$\cardeq{X}{Y}$ అని సంక్షిప్తంగా రాస్తాం. (సమితులు $X$, $Y$ గురించి
అనధికారికంగా మాట్లాడేటప్పుడు కూడా ఇదే సంజ్ఞామానాన్ని వాడతాం కాబట్టి
ఇది కొంత గందరగోళంగా ఉండవచ్చు---కానీ ఇక్కడ ఈ సంజ్ఞామానం ఏకస్థాన
సంబంధ చరాలు $X$,~$Y$ ఉన్న ద్వితీయ-స్థాయి తర్క
\tetoken{సూత్రాలకు}{!!{formula}s} సంక్షిప్త రూపం.)

\begin{prop}
\tetoken{వాక్యం}{!!{sentence}} $\lforall[X][\lforall[Y][((\cardle{X}{Y} \land
    \cardle{Y}{X}) \lif \cardeq{X}{Y})]]$ చెల్లుబాటవుతుంది.
\end{prop}

\begin{proof}
ఏ ఉపసమితులు $X \subseteq \Domain{M}$, $Y \subseteq \Domain{M}$కైనా
$\cardle{X}{Y}$, $\cardle{Y}{X}$ అయితే $\cardeq{X}{Y}$ అయ్యే
పక్షంలో, \tetoken{వాక్యం}{!!{sentence}} \tetoken{ఒక నిర్మాణం}{!!a{structure}}~$\Struct{M}$లో
సంతృప్తిపరచబడుతుంది. కానీ ఇది \emph{ఏ} సమితులు $X$, $Y$కైనా
వర్తిస్తుంది---ఇదే ష్రోడర్--బెర్న్‌స్టైన్ సిద్ధాంతం.
\end{proof}


\end{document}
